% Part: sets-functions-relations
% Chapter: arithmetization
% Section: rationals

\documentclass[../../../include/open-logic-section]{subfiles}

\begin{document}

\olfileid{sfr}{arith}{rat}
\olsection{له $\Int$ څخه $\Rat$ ته}

همدا اوس مو وليدل چې په ساده سټ نظريه څنګه له طبيعي عددونو څخه
صحيح عددونه جوړېږي۔ اوس به په ډېرې ورته لارې له صحيح عددونو څخه
ناطق عددونه جوړول وګورو۔ زموږ لومړنۍ خبرې دا دي:
\begin{enumerate}
	\item هر ناطق عدد د $\nicefrac{i}{j}$ په بڼه ليکل کېدے شي،
	چې $i$ او $j$ دواړه صحيح عددونه وي خو $j$ صفر نۀ وي۔
	\item په عبارت $\nicefrac{i}{j}$ کښې زېرمه شوي معلومات په هماغه
	ډول په مرتبه جوړه $\tuple{ i, j}$ کښې زېرمه کېدے شي۔
\end{enumerate}
ښکاره لار به دا وي چې ناطق عددونه د $\Int \times
(\Int \setminus \{0_\Int\})$ څخه د اخيستل شوو مرتبو جوړو \emph{په
توګه} وګڼو۔ خو لکه مخکښې، دا به لږ ډېر ساده وي، ځکه موږ غواړو
$\nicefrac{3}{2} = \nicefrac{6}{4}$ وي، خو $\tuple{ 3, 2}\neq
\tuple{ 6, 4}$۔ په عمومي ډول لاندې خاصيت غواړو:
\[
	\nicefrac{a}{b} = \nicefrac{c}{d} \text{ هله او يوازې هله چې } a \times d = b \times c
\]
د دې د ترلاسه کولو دپاره پر $\Int \times (\Int \setminus
\{0_\Int\})$ د هم ارزښتۍ يوه اړيکه داسې تعريفوو:
\[
	\tuple{ a, b }\Ratequiv \tuple{ c, d} \text{ هله او يوازې هله چې }a \times d = b \times c
\]
بايد وارزوو چې دا د هم ارزښتۍ اړيکه ده۔ دا د $\Intequiv$ حالت ته
ډېر ورته دے، او موږ ئې د تمرين په توګه پرېږدو۔
\begin{prob}
وښيئ چې $\Ratequiv$ د هم ارزښتۍ اړيکه ده۔
\end{prob}
خو دا موږ ته د لاندې خبرې اجازه راکوي:
\begin{defn}
ناطق عددونه د $\Ratequiv$ له مخې د صحيح عددونو د جوړو د هم ارزښتۍ
ټولګي دي (چې دويم غړے ئې صفر نۀ وي)۔ يعنې $\Rat =
\equivclass{(\Int \times (\Int\setminus \{0_\Int\}))}{\Ratequiv}$۔
\end{defn}

لکه د صحيح عددونو په شان، ځينې بنسټيزې عمليې هم تعريفول غواړو۔
چې $\equivrep{i,j}{\Ratequiv}$ د $\Ratequiv$ له مخې هغه د هم
ارزښتۍ ټولګے وي چې $\tuple{i, j}$ ئې يو !!a{element} دے، وايو:
\begin{align*}
	\equivrep{a, b}{\Ratequiv} + \equivrep{c, d}{\Ratequiv} &= \equivrep{ad + bc,  bd}{\Ratequiv}\\
	\equivrep{a, b}{\Ratequiv} \times \equivrep{c, d}{\Ratequiv} &= \equivrep{a  c, b d}{\Ratequiv}.
\intertext{پر دې ناطقو عددونو د $r \leq s$ د تعريفولو دپاره دا حقيقت کاروو چې
$r \le s$ هله او يوازې هله چې $s - r$ منفي نۀ وي؛ يعنې $r - s$ د
$\nicefrac{i}{j}$ په بڼه ليکل کېدے شي چې $i$ نامنفي او $j$~مثبت وي۔
\textbf{د ژباړې د نسخې سمون (OLARI-001):} په انګرېزي نثر کښې د وروستي توپير ترتيب سرچپه ليکل شوے؛ لاندې معادله د نامنفي توپير سم ترتيب ښيي:}
	\equivrep{a, b}{\Ratequiv} \leq \equivrep{c, d}{\Ratequiv} &\text{ هله او يوازې هله چې }
	\equivrep{c, d}{\Ratequiv} - \equivrep{a, b}{\Ratequiv} =
	\equivrep{i_\Int, j_\Int}{\Ratequiv}
\end{align*}
د ځينو $i \in \Nat$ او $0 \neq j \in \Nat$ دپاره۔

بيا بايد وارزوو چې دا تعريفونه لکه څنګه چې \emph{بايد}، هماغسې
چلېږي؛ او دا کار \olref[check]{sec} ته پرېږدو۔ خو رښتيا هم سم
چلېږي!{} په پاے کښې صحيح عددونه د ناطقو عددونو \emph{په توګه}
د کارولو يوه لار غواړو؛ نو د هر $i \in \Int$ دپاره ټاکو چې
$i_\Rat = \equivrep{i, 1_\Int}{\Ratequiv}$۔ يو ځل بيا په
\olref[check]{sec} کښې ارزوو چې دا ټول سم چلند کوي۔

\begin{prob}
وښيئ چې د هر $i, j \in \Int$ دپاره $(i + j)_\Rat = i_\Rat+ j_\Rat$
او $(i \times j)_\Rat = i_\Rat \times j_\Rat$ او
$i \leq j \liff i_\Rat \leq j_\Rat$۔
\end{prob}

\end{document}
